\documentclass[12pt]{article}
\usepackage{amsmath, amsfonts, amssymb, amsthm, mathdots, mathrsfs, dsfont,systeme}
\usepackage[utf8]{inputenc}
\usepackage[spanish]{babel}
\usepackage{color}
\usepackage{fancyhdr}
\usepackage{enumerate}
\usepackage{hyperref}
\hypersetup{
    colorlinks=true,
    linkcolor=blue,
    filecolor=magenta,   
    citecolor=blue,
    urlcolor=cyan,
    pdftitle={Algebra lineal},
    pdfpagemode=FullScreen,
    }
\usepackage{amscd}
\usepackage{nomencl}
\usepackage{tikz}
\usepackage{tikz-cd}
\usepackage{amssymb}
\usepackage{blkarray}
\usepackage{multicol}
\usepackage[a4paper,margin=1cm,footskip=.5cm]{geometry}
\usepackage[all]{xy}
\usepackage{pst-node}
\usepackage{subfigure}
%\usepackage[table]{xcolor}
\usepackage{colortbl}
\newcommand{\cN}{\mathcal{N}}
\newcommand{\V}{\mathbb{V}}
\newcommand{\A}{\mathbb{A}}
\newcommand{\M}{\mathcal{M}}
\newcommand{\cH}{\mathcal{H}}
\newcommand{\F}{\mathcal{F}}
\newcommand{\G}{\mathcal{G}}
\newcommand{\cL}{\mathcal{L}}
\newcommand{\cO}{\mathcal{O}}
\newcommand{\Rn}{\mathbb{R}^n}
\DeclareMathOperator{\Hom}{Hom}
\DeclareMathOperator{\Ns}{NS}
\DeclareMathOperator{\im}{im}
\DeclareMathOperator{\Pic}{Pic}
\DeclareMathOperator{\Div}{Div}
\DeclareMathOperator{\Pdiv}{PDiv}
\DeclareMathOperator{\ddiv}{div}
\DeclareMathOperator{\Cl}{Cl}
\DeclareMathOperator{\CaDiv}{CaDiv}
\DeclareMathOperator{\Mor}{Mor}
\DeclareMathOperator{\Spec}{Spec}
\DeclareMathOperator{\id}{id}
\DeclareMathOperator{\Aut}{Aut}
\DeclareMathOperator{\End}{End}
\DeclareMathOperator{\CH}{CH}
\DeclareMathOperator{\IM}{Im}
\DeclareMathOperator{\Pf}{Pf}
\DeclareMathOperator{\Gl}{GL}
\DeclareMathOperator{\Sl}{SL}
\DeclareMathOperator{\St}{St}
\DeclareMathOperator{\tr}{Tr}
\DeclareMathOperator{\Gr}{Gr}
\DeclareMathOperator{\diag}{diag}
\DeclareMathOperator{\ord}{ord}
\DeclareMathOperator{\supp}{supp}
\DeclareMathOperator{\sign}{sign}
\DeclareMathOperator{\adj}{adj}
\DeclareMathOperator{\nul}{nul}
\DeclareMathOperator{\ran}{ran}
\newcommand{\calP}{\cal{P}}
\newcommand{\p}{\mathbb{P}}
\newcommand{\R}{\mathbb{R}}
\newcommand{\Z}{\mathbb{Z}}
\newcommand{\Q}{\mathbb{Q}}
\newcommand{\C}{\mathbb{C}}
\newcommand{\N}{\mathbb{N}}
\newcommand{\bbH}{\mathbb{H}}
\setlength{\parindent}{0pt}
\setlength{\textwidth}{6.5in}
\setlength{\textheight}{9.0in}
\setlength{\topmargin}{0in}
\setlength{\oddsidemargin}{0in}
\newtheoremstyle{exampstyle}
  {\topsep} % Space above
  {\topsep} % Space below
  {\itshape} % Body font
  {} % Indent amount
  {\bfseries} % Theorem head font
  {.} % Punctuation after theorem head
  {.5em} % Space after theorem head
  {} % Theorem head spec (can be left empty, meaning `normal')
\theoremstyle{exampstyle}
\newtheorem{thm}{Teorema}[section]
\newtheorem{prop}[thm]{Proposici\'on}
\newtheorem{lemma}[thm]{Lemma}
\newtheorem{cor}[thm]{Corolario}
\theoremstyle{definition}
\newtheorem{thme}{Expected Theorem}

\newtheorem{rec}[thm]{Recuerdo}
\newtheorem{conj}{Conjecture}
\newtheorem{defin}[thm]{Definici\'on}

\newtheoremstyle{break}%
{}{}%
{}{}%
{\bfseries}{}% % Note that final punctuation is omitted.
{\newline}{}
\theoremstyle{break}
\newtheorem{ejem}[thm]{Ejemplo}
\newtheorem{ejer}[thm]{Ejercicio}
\newtheorem{sol}[thm]{Soluci\'on}
\newtheorem{rmk}[thm]{Observaci\'on}
\numberwithin{equation}{section}
\newcommand{\gb}[1]{\operatorname{\textbf {#1}}}
\hypersetup{linkcolor=black}
\title{\'Algebra Lineal}
\author{Prof. Sebasti\'an Cea \\Prof. Camila Muñoz\\
}
\date{}
\pagestyle{fancy}

\lhead{\includegraphics[width = 4cm]{UFT.png}}
\setlength{\headsep}{0.6in}
\newcolumntype{a}{ >{\columncolor{yellow}} c }

\fancyhead[RO,LE]{
Prof: C. Muñoz.\\
02/2023}}



\begin{document}
%\maketitle
%\newpage

\begin{minipage}[r]{15cm}


\end{minipage}

\begin{center}
\Large{\textbf{\'Algebra Lineal\\Taller evaluado 3  }}
\end{center}

\section{Ejercicios}


\begin{enumerate}
    \item Considere el conjunto $S=\{ \begin{pmatrix} 0 \\ 1 \\ 2 \end{pmatrix}, \begin{pmatrix} 1 \\ 1 \\ 2 \end{pmatrix}, \begin{pmatrix} 0 \\ 0 \\ -1 \end{pmatrix} \}$.
    \begin{enumerate}
        \item Muestre que $S$ es base de $\R^{3}$.
        \item Encuentre una base ortonormal de $\R^{3}$ a partir de esta base usando el procedimiento de Gramm-Schmidt.
        \item Encuentre las coordenadas de $\{ \begin{pmatrix} -1 \\0 \\ 1 \end{pmatrix} \}$ en esta base.    
    \item Encuentre la proyección de $\{ \begin{pmatrix} 1 \\ 0\\ 1 \end{pmatrix} \}$ sobre el plano generado por $\{ \begin{pmatrix} 0 \\ 1 \\ 2 \end{pmatrix}, \begin{pmatrix} 1 \\ 1 \\ 2 \end{pmatrix} \}$.

\end{enumerate} 
     
\item  Considere la función $T: \R^{} \to \R^{3}$ dada por $T \begin{pmatrix} x \\ y \\ z \end{pmatrix} = \begin{pmatrix}
    x+z \\ y+z \\ 2z 
\end{pmatrix}$

\begin{enumerate}
\item  Demuestre que $T$ es lineal.
    \item Determine $\ker T$ y $\dim \operatorname{Im} T$.
    \item Encuentre su matriz $[T]_C$ representante en la base canónica $C$. Luego, usando dicha matriz, eval\'ue $T(1,2,3)$.
    \item Determine el polinomio caracter\'istico de $[T]_C$, encuentre sus valores propios y determine sus subespacios propios. De ser posible, encuentre la matriz $P$, tal que $[T]_C=PDP^{-1}$, con $D$ matriz diagonal. En caso negativo, justifique.
    \end{enumerate}

 

\section*{Instrucciones}
\begin{itemize}
    \item Grupos de entre 4 y 5 personas.
    \item Se entrega el d\'ia jueves 30 a las 23:59.
\end{itemize}
\end{document}